\section{Week 2: Anatomy of a Model: Signal Detection and Forgetting
Curves}

\begin{quote}
\textbf{Chapter 1: Fundamentals of cognitive modeling} · Lecture 1 h ·
Lab 2 h \textbf{This week's running models:} equal-variance Gaussian
signal detection theory (SDT); exponential and power forgetting curves.
Both return in Weeks 3--13.
\end{quote}

\begin{center}\rule{0.5\linewidth}{0.5pt}\end{center}

\subsection{Key terms}

Important terms of this lesson, with their Persian equivalents (repeated
even if they appeared in earlier weeks).

{\def\LTcaptype{none} % do not increment counter
\begin{longtable}[]{@{}
  >{\raggedright\arraybackslash}p{(\linewidth - 4\tabcolsep) * \real{0.3333}}
  >{\raggedright\arraybackslash}p{(\linewidth - 4\tabcolsep) * \real{0.3333}}
  >{\raggedright\arraybackslash}p{(\linewidth - 4\tabcolsep) * \real{0.3333}}@{}}
\toprule\noalign{}
\begin{minipage}[b]{\linewidth}\raggedright
English
\end{minipage} & \begin{minipage}[b]{\linewidth}\raggedright
\RL{\CMfa فارسی}
\end{minipage} & \begin{minipage}[b]{\linewidth}\raggedright
\RL{\CMfa توضیح کوتاه}
\end{minipage} \\
\midrule\noalign{}
\endhead
\bottomrule\noalign{}
\endlastfoot
forgetting curve & \RL{\CMfa منحنی فراموشی} & \RL{\CMfa کاهش یادآوری با گذشت
زمان} \\
signal detection theory (SDT) & \RL{\CMfa نظریه‌ی تشخیص علامت} & \RL{\CMfa مدلی که
حساسیت را از سوگیری پاسخ جدا می‌کند} \\
exponential / power function & \RL{\CMfa تابع نمایی} / \RL{\CMfa توانی} & \RL{\CMfa دو شکل
رایج برای منحنی فراموشی و یادگیری} \\
parameter & \RL{\CMfa پارامتر} & \RL{\CMfa کمیت قابل تنظیم مدل که بین افراد یا شرایط
فرق می‌کند} \\
deterministic core & \RL{\CMfa هسته‌ی قطعی} & \RL{\CMfa بخش بدون تصادف مدل که پیش‌بینی
میانگین را می‌دهد} \\
noise & \RL{\CMfa نویز} & \RL{\CMfa تغییرات تصادفی در شواهد یا در داده} \\
linking assumption & \RL{\CMfa فرض پیوند} & \RL{\CMfa اینکه کمیت‌های مدل چطور به
رفتار قابل مشاهده وصل می‌شوند} \\
sensitivity (d′) & \RL{\CMfa حساسیت} & \RL{\CMfa فاصله‌ی توزیع علامت و نویز} \\
hit & \RL{\CMfa اصابت} (hit) & \RL{\CMfa پاسخ «بله» وقتی علامت وجود دارد} \\
response bias & \RL{\CMfa سوگیری پاسخ} & \RL{\CMfa تمایل به یک پاسخ، مستقل از
محرک} \\
binomial distribution & \RL{\CMfa توزیع دوجمله‌ای} & \RL{\CMfa توزیع تعداد موفقیت در}
n \RL{\CMfa کوشش} \\
noise model & \RL{\CMfa مدل نویز} & \RL{\CMfa توزیع داده حول پیش‌بینی مدل} \\
simulation & \RL{\CMfa شبیه‌سازی} & \RL{\CMfa تولید داده‌ی مصنوعی با اجرای مدل} \\
ROC / zROC curve & \RL{\CMfa منحنی} ROC / zROC & \RL{\CMfa نرخ} hit \RL{\CMfa در برابر
نرخ} false alarm \RL{\CMfa با تغییر ملاک} \\
fit / fitting & \RL{\CMfa برازش} & \RL{\CMfa تنظیم پارامترها تا مدل با داده جور
شود} \\
criterion (c) & \RL{\CMfa ملاک} & \RL{\CMfa آستانه‌ی تصمیم؛ نشان‌دهنده‌ی سوگیری
پاسخ} \\
formal model / formalize & \RL{\CMfa مدل فرمال} / \RL{\CMfa فرمال کردن} & \RL{\CMfa نوشتن
نظریه به شکل معادله یا الگوریتم} \\
trial & \RL{\CMfa کوشش} (trial) & \RL{\CMfa یک بار انجام تکلیف توسط شرکت‌کننده} \\
false alarm & \RL{\CMfa هشدار کاذب} (false alarm) & \RL{\CMfa پاسخ «بله» وقتی علامت
وجود ندارد} \\
estimate / estimator & \RL{\CMfa برآورد} / \RL{\CMfa برآوردگر} & \RL{\CMfa مقدار به‌دست‌آمده
برای پارامتر} / \RL{\CMfa روش به دست آوردن آن} \\
miss / correct rejection & \RL{\CMfa از دست دادن} / \RL{\CMfa رد درست} & \RL{\CMfa پاسخ
«نه» با وجود علامت} / \RL{\CMfa پاسخ «نه» بدون علامت} \\
liberal / conservative & \RL{\CMfa آسان‌گیر} / \RL{\CMfa محافظه‌کار} & \RL{\CMfa تمایل به
گفتن «بله»} / \RL{\CMfa «نه»} \\
free parameter & \RL{\CMfa پارامتر آزاد} & \RL{\CMfa پارامتری که از روی داده برآورد
می‌شود} \\
verbal theory & \RL{\CMfa نظریه‌ی کلامی} & \RL{\CMfa نظریه‌ای که فقط با کلمات بیان شده
و پیش‌بینی دقیق نمی‌دهد} \\
\end{longtable}
}

\begin{center}\rule{0.5\linewidth}{0.5pt}\end{center}

\subsection{Learning outcomes}

By the end of this week, students will be able to:

\begin{enumerate}
\def\labelenumi{\arabic{enumi}.}
\tightlist
\item
  \textbf{Decompose} a cognitive model into its components: design, free
  and fixed parameters, deterministic core, noise (data) model, and
  linking assumptions.
\item
  \textbf{Compute} \(d'\) and \(c\) from a 2×2 table of counts,
  including a correction for hit or false-alarm rates of 0 or 1, and
  \textbf{explain} why \(d'\) separates sensitivity from bias.
\item
  \textbf{Implement} exponential and power forgetting models with a
  binomial noise model, and \textbf{predict} qualitatively how each
  parameter changes the curves.
\item
  \textbf{Use} simulation and diagnostic plots (ROC/zROC, semi-log and
  log-log plots) to \textbf{evaluate} whether a model's qualitative
  signature appears in data.
\item
  \textbf{Interpret} fitted parameters in psychological terms and
  \textbf{identify} the limits of the ``fit a straight line to
  transformed data'' approach.
\end{enumerate}

\begin{center}\rule{0.5\linewidth}{0.5pt}\end{center}

\subsection{Motivation \& big picture}

\textbf{Two radiologists.} Radiologist R1 flags 93\% of scans that
really contain a tumor, but also flags 31\% of healthy scans.
Radiologist R2 flags only 69\% of tumors, but just 7\% of healthy scans.
Their percent correct is identical (81\%). Who is the better
diagnostician?

Percent correct cannot answer this, because it mixes two different
things: \textbf{how well} a person can tell tumors from healthy tissue
(\emph{sensitivity}), and \textbf{how willing} they are to say ``tumor''
(\emph{bias}). Signal detection theory is a small model with two
parameters that pulls these apart. As the worked example will show, R1
and R2 have essentially the \textbf{same} sensitivity; they differ only
in where they set their criterion. This is the first payoff of modeling:
model parameters can \textbf{measure} latent constructs that raw scores
conflate.

\textbf{Ebbinghaus's curve.} Since Ebbinghaus's experiments on himself
in the 1880s, we know that forgetting is fast at first and then slows
down. ``Fast, then slow'' is a verbal description that both an
\textbf{exponential} and a \textbf{power} function satisfy. Yet they
imply different things: an exponential means a memory loses a
\emph{constant proportion} of its strength per day, at any age; a power
function means older memories become \emph{more resistant}. Choosing
between them is a question about the nature of memory, and it requires a
formal model.

\paragraph{Where is this used?}

\begin{itemize}
\tightlist
\item
  \textbf{Medical diagnosis.} Radiologists, pathologists and screening
  programs are evaluated with SDT and ROC curves, so that a doctor who
  ``sees better'' can be told apart from a doctor who simply says
  ``abnormal'' more often.
\item
  \textbf{Memory research.} In a recognition test, studied words are
  ``signal'' and new words are ``noise''. \(d'\) measures memory
  strength separately from the tendency to say ``old''.
\item
  \textbf{Eyewitness identification.} Police line-up procedures are
  compared with ROC analysis: some procedures do not improve accuracy,
  they only make witnesses more cautious.
\item
  \textbf{Perception and hearing tests.} Thresholds for seeing or
  hearing faint stimuli are measured with SDT, so that a cautious
  listener is not mistaken for a hard-of-hearing one.
\item
  \textbf{Machine learning.} Every binary classifier (spam filter, fraud
  detector, AI-text detector) has a hit rate and a false-alarm rate, and
  is evaluated with ROC curves and the area under the curve (AUC). This
  is the same theory.
\item
  \textbf{Forgetting curves} are used to schedule learning:
  spaced-repetition apps (e.g.~for vocabulary) decide when to show an
  item again from a model of how fast it is forgotten.
\end{itemize}

\begin{center}\rule{0.5\linewidth}{0.5pt}\end{center}

\subsection{Lecture notes}

\paragraph{In plain words}

This week has one big idea: a model has \textbf{parts}, and each part is
a separate assumption. We look at two small models. \textbf{Signal
detection theory} explains yes/no decisions under uncertainty. Every
trial gives the observer a noisy ``feeling of signal''. If the feeling
is stronger than a personal threshold, the observer says ``yes''. Two
numbers describe the observer: how far apart the ``signal'' and ``no
signal'' feelings are on average (\textbf{sensitivity}, \(d'\)), and
where the threshold is (\textbf{criterion}, \(c\)). Two people can have
the same percent correct for completely different reasons, and SDT shows
which. The \textbf{forgetting curve} describes how recall drops with
time. Two formulas, exponential and power, both look like ``fast drop,
then slow'', but they say different things about memory. Choosing
between them is the kind of question the rest of the course teaches us
to answer with data.

\paragraph{Statistics words in plain language}

\begin{itemize}
\tightlist
\item
  \textbf{Distribution:} a description of which values a random quantity
  takes and how often. A histogram of many measurements approximates it.
\item
  \textbf{Normal (Gaussian) distribution \(\mathcal{N}(\mu, \sigma)\):}
  the bell curve, with mean \(\mu\) (center) and standard deviation
  \(\sigma\) (width). About 68\% of values fall within
  \(\mu \pm \sigma\) and 95\% within \(\mu \pm 1.96\sigma\).
\item
  \textbf{Standard normal:} the normal distribution with \(\mu = 0\),
  \(\sigma = 1\). A value measured on this scale is called a
  \textbf{z-score}: ``how many SDs above the mean''.
\item
  \textbf{\(\Phi(x)\) (CDF):} the probability that a standard normal
  value is \emph{below} \(x\). \(\Phi(0) = 0.5\),
  \(\Phi(1) \approx 0.84\), \(\Phi(1.96) \approx 0.975\). In Python:
  \nolinkurl{norm.cdf(x)}.
\item
  \textbf{\(\Phi^{-1}(p)\) (inverse CDF, ``z-transform''):} the opposite
  question: below which value lies a proportion \(p\)?
  \(\Phi^{-1}(0.84) \approx 1\). In Python: \nolinkurl{norm.ppf(p)}.
\item
  \textbf{Binomial distribution \(\mathrm{Binomial}(n, p)\):} the number
  of successes in \(n\) independent tries, each with success probability
  \(p\) (e.g.~words recalled out of 50). Its mean is \(np\) and its
  variance \(np(1-p)\).
\item
  \textbf{Proportion vs.~probability:} a proportion is what we
  \emph{observe} (30 of 50 recalled = 0.60); a probability is the
  model's \emph{underlying} value that the proportion estimates.
\end{itemize}

\textbf{Suggested timing (60 min):} anatomy of a model 15 · SDT 20 ·
forgetting curves 15 · simulation as a thinking tool + preview 10.

\subsubsection{1. The anatomy of a model}

In Week 1 we wrote a model as a data-generating machine,
\(\mathbf{y} \sim p(\mathbf{y}\mid\boldsymbol\theta, \mathcal{M})\). In
practice, almost every model you will meet in this course can be taken
apart into five pieces:

{\def\LTcaptype{none} % do not increment counter
\begin{longtable}[]{@{}
  >{\raggedright\arraybackslash}p{(\linewidth - 6\tabcolsep) * \real{0.2500}}
  >{\raggedright\arraybackslash}p{(\linewidth - 6\tabcolsep) * \real{0.2500}}
  >{\raggedright\arraybackslash}p{(\linewidth - 6\tabcolsep) * \real{0.2500}}
  >{\raggedright\arraybackslash}p{(\linewidth - 6\tabcolsep) * \real{0.2500}}@{}}
\toprule\noalign{}
\begin{minipage}[b]{\linewidth}\raggedright
Component
\end{minipage} & \begin{minipage}[b]{\linewidth}\raggedright
Question
\end{minipage} & \begin{minipage}[b]{\linewidth}\raggedright
SDT
\end{minipage} & \begin{minipage}[b]{\linewidth}\raggedright
Forgetting (exponential)
\end{minipage} \\
\midrule\noalign{}
\endhead
\bottomrule\noalign{}
\endlastfoot
\textbf{1. Design} \(x_i\) & What is manipulated and measured? & Trial
type (signal/noise); response yes/no & Retention interval \(t_i\);
number of items recalled out of \(n\) \\
\textbf{2. Parameters} \(\boldsymbol\theta\) & What is free to vary?
What is fixed? & Free: \(d'\), \(c\). Fixed: both SDs = 1, noise mean =
0 & Free: \(A\), \(\lambda\). Fixed: \(n\) \\
\textbf{3. Deterministic core} \(f(x;\boldsymbol\theta)\) & What does
the model predict on average? & \(H = \Phi(d'/2 - c)\),
\(F = \Phi(-d'/2 - c)\) & \(f(t) = A e^{-\lambda t}\) \\
\textbf{4. Noise model} & How do observations scatter around the
prediction? & Hits \(\sim \mathrm{Binomial}(n_s, H)\); false alarms
\(\sim \mathrm{Binomial}(n_n, F)\) & Recalled
\(\sim \mathrm{Binomial}(n, f(t))\) \\
\textbf{5. Linking assumptions} & How do model quantities map onto
behavior? & ``Yes'' if internal evidence exceeds the criterion & Recall
probability equals memory strength \(f(t)\) \\
\end{longtable}
}

Some points about each component:

\begin{itemize}
\tightlist
\item
  \textbf{Free vs.~fixed parameters.} Fixed parameters are not
  ``unimportant''. They are \emph{assumptions}. In Week 1 we fixed
  \(t_0 = 300\) ms. In SDT, fixing the noise mean at 0 and both standard
  deviations at 1 is what makes \(d'\) and \(c\) estimable at all:
  shifting or stretching the whole axis would not change any prediction.
  Such fixes are called \textbf{scaling constraints}.
\item
  \textbf{Parameter ranges.} Every parameter has a meaningful range:
  \(0 \le A \le 1\), \(\lambda > 0\), \(d' \ge 0\) for a non-perverse
  observer. Respecting these ranges matters when fitting (Week 3).
\item
  \textbf{The noise model is part of the model.} Two researchers with
  the same \(f(t)\) but different noise models (binomial vs.~Gaussian)
  have \emph{different} models and will get different parameter
  estimates (Week 4). The binomial is natural for counts: its variance
  \(f(1-f)/n\) is largest at \(f = 0.5\) and shrinks near 0 and 1.
\end{itemize}

\subsubsection{2. Signal detection theory (SDT)}

\textbf{The task.} On each trial a stimulus is either \emph{signal}
(e.g., a tumor, a studied word, a faint tone) or \emph{noise}. The
observer answers ``yes'' (signal) or ``no''. The outcomes form a 2×2
table:

{\def\LTcaptype{none} % do not increment counter
\begin{longtable}[]{@{}lll@{}}
\toprule\noalign{}
& Respond ``yes'' & Respond ``no'' \\
\midrule\noalign{}
\endhead
\bottomrule\noalign{}
\endlastfoot
\textbf{Signal trial} & Hit & Miss \\
\textbf{Noise trial} & False alarm (FA) & Correct rejection (CR) \\
\end{longtable}
}

with hit rate \(H = \frac{\text{hits}}{\text{hits} + \text{misses}}\)
and false-alarm rate \(F = \frac{\text{FAs}}{\text{FAs} + \text{CRs}}\).

\textbf{The model (intuition).} Each trial produces an internal
\textbf{evidence} value. It is noisy: even pure noise trials sometimes
feel signal-like. Signal trials produce, on average, more evidence. The
observer says ``yes'' whenever the evidence exceeds a
\textbf{criterion}.

\textbf{The model (formal), equal-variance Gaussian SDT:}

\[
\text{noise trials: } X \sim \mathcal{N}(0, 1), \qquad
\text{signal trials: } X \sim \mathcal{N}(d', 1), \qquad
\text{respond "yes" iff } X > k.
\]

\begin{itemize}
\tightlist
\item
  \(d'\) is \textbf{sensitivity}: the distance between the two
  distributions in SD units.
\item
  \(k\) is the absolute criterion location. It is more convenient to
  measure the criterion relative to the midpoint between the
  distributions: \(c = k - d'/2\). With \(c = 0\) the observer is
  unbiased, \(c < 0\) is \textbf{liberal} (says ``yes'' readily), and
  \(c > 0\) is \textbf{conservative}.
\end{itemize}

\textbf{Predictions.} Using the standard normal CDF \(\Phi\):

\[
H = P(X > k \mid \text{signal}) = \Phi\!\left(\tfrac{d'}{2} - c\right), \qquad
F = P(X > k \mid \text{noise}) = \Phi\!\left(-\tfrac{d'}{2} - c\right).
\]

\textbf{Inverting the model.} Applying \(\Phi^{-1}\) (often written
\(z\)) to both equations and solving gives the familiar formulas:

\[
d' = \Phi^{-1}(H) - \Phi^{-1}(F), \qquad
c = -\tfrac{1}{2}\left[\Phi^{-1}(H) + \Phi^{-1}(F)\right].
\]

Notice what just happened: this is our first \textbf{parameter
estimation}. There are two free parameters and two data points (\(H\),
\(F\)), so the model can be solved exactly. With more data points than
parameters, an exact solution usually does not exist and we need the
fitting methods of Weeks 3--4.

\textbf{The extreme-rate problem.} If \(H = 1\) or \(F = 0\), then
\(\Phi^{-1}\) is \(\pm\infty\) and \(d'\) is infinite. This happens
often with few trials. A standard fix is the \textbf{log-linear
correction}: add 0.5 to every cell count before computing the rates
(i.e., add 0.5 to hits and FAs and 1 to the number of trials of each
type). We use it by default in the lab.

\textbf{The ROC curve: SDT's signature.} Hold \(d'\) fixed and let \(c\)
vary. The pair \((F, H)\) traces a curve, the \textbf{receiver operating
characteristic (ROC)}. Every point on one ROC has the same sensitivity
but a different bias. Plotting \(\Phi^{-1}(H)\) against \(\Phi^{-1}(F)\)
gives the \textbf{zROC}, which under equal-variance SDT is a straight
line:

\[
\Phi^{-1}(H) = d' + 1 \cdot \Phi^{-1}(F).
\]

So the model makes a \emph{testable} prediction: \textbf{the zROC is
linear with slope 1}. ROCs can be measured with confidence ratings (each
confidence level acts as a separate criterion). In recognition memory,
zROC slopes are typically \emph{below} 1, which motivates the
unequal-variance model below.

\begin{quote}
\textbf{Advanced: unequal-variance SDT.} If signal trials have
\(X \sim \mathcal{N}(\mu_s, \sigma_s)\), the zROC is still linear but
its slope is \(1/\sigma_s\). In that case, \(d'\) computed from a single
\((H, F)\) pair \emph{changes with the criterion}, so it is no longer a
pure sensitivity measure. A common alternative index is
\(d_a = \mu_s \sqrt{2/(1+\sigma_s^2)}\). Estimating \(\sigma_s\)
requires more than one criterion, i.e., rating data. You will see this
effect directly in lab Part B.
\end{quote}

\paragraph{Reading the math}

\textbf{The two bell curves.} Draw two normal curves on one axis: the
``noise'' curve centered at 0 and the ``signal'' curve centered at
\(d'\). Both have width 1. Draw a vertical line at the criterion \(k\).
Everything to the right of the line is a ``yes''.

\begin{itemize}
\tightlist
\item
  \(H\) = area of the \textbf{signal} curve to the right of the line.
  \(F\) = area of the \textbf{noise} curve to the right of the line.
\item
  Why \(H = \Phi(d'/2 - c)\)? The criterion sits at \(k = d'/2 + c\)
  (the midpoint plus \(c\)). The signal curve is centered at \(d'\), so
  the line is \(k - d' = c - d'/2\) SDs from its center. The area to the
  right of that point is \(1 - \Phi(c - d'/2) = \Phi(d'/2 - c)\),
  because the normal curve is symmetric.
\item
  \textbf{Numerical example.} \(d' = 2\), \(c = 0\):
  \(H = \Phi(1) = 0.84\) and \(F = \Phi(-1) = 0.16\). Now make the
  observer liberal, \(c = -0.5\): \(H = \Phi(1.5) = 0.93\),
  \(F = \Phi(-0.5) = 0.31\). Both rates went up, sensitivity is
  unchanged. This is radiologist R1 of the worked example.
\end{itemize}

\textbf{Inverting.} \(d' = \Phi^{-1}(H) - \Phi^{-1}(F)\): convert each
rate into a z-score and subtract. For R1: \(1.476 - (-0.496) = 1.97\).
\(c\) is minus the average of the two z-scores: if both rates are high
(both z-scores positive), \(c\) is negative, i.e.~liberal.

\textbf{The zROC line.} \(\Phi^{-1}(H) = d' + \Phi^{-1}(F)\) is just the
previous equation rearranged. In z-units, moving the criterion moves
both z-scores by the same amount, so the points lie on a line with slope
1 and intercept \(d'\).

\paragraph{For an engineer}

SDT \emph{is} the detection theory of radar and communications, imported
into psychology in the 1950s.

{\def\LTcaptype{none} % do not increment counter
\begin{longtable}[]{@{}
  >{\raggedright\arraybackslash}p{(\linewidth - 2\tabcolsep) * \real{0.5000}}
  >{\raggedright\arraybackslash}p{(\linewidth - 2\tabcolsep) * \real{0.5000}}@{}}
\toprule\noalign{}
\begin{minipage}[b]{\linewidth}\raggedright
Psychology
\end{minipage} & \begin{minipage}[b]{\linewidth}\raggedright
Telecom / radar
\end{minipage} \\
\midrule\noalign{}
\endhead
\bottomrule\noalign{}
\endlastfoot
noise trial / signal trial & hypothesis \(H_0\) (noise only) / \(H_1\)
(signal + noise) \\
internal evidence \(X\) & test statistic, e.g.~matched-filter output \\
criterion \(k\) (or \(c\)) & detection threshold \\
hit rate \(H\) & probability of detection \(P_D\) \\
false-alarm rate \(F\) & probability of false alarm \(P_{FA}\) \\
sensitivity \(d'\) & deflection \((\mu_1 - \mu_0)/\sigma\); \(d'^2\)
plays the role of an SNR \\
ROC curve & receiver operating characteristic (same name, same plot) \\
liberal / conservative observer & low / high threshold (high / low
\(P_{FA}\)) \\
\end{longtable}
}

Differences in emphasis: an engineer usually \emph{chooses} the
threshold (e.g.~Neyman--Pearson: fix \(P_{FA}\), maximize \(P_D\)). A
psychologist \emph{observes} a person's \(P_D\) and \(P_{FA}\) and
\emph{infers} both the SNR (\(d'\)) and the threshold (\(c\)) the person
used. The criterion \(c\) is measured from the midpoint between the two
means, so \(c = 0\) means ``no bias''. The \textbf{optimal} threshold
depends on prior probabilities and costs exactly as in Bayesian
detection, and people often move their criterion in the right direction
when payoffs change (Week 4).

\subsubsection{3. Forgetting curves}

\textbf{Data.} At each retention interval \(t_i\) (e.g., in days), the
participant is tested on \(n\) items and recalls \(y_i\) of them.

\textbf{Two candidate deterministic cores:}

{\def\LTcaptype{none} % do not increment counter
\begin{longtable}[]{@{}
  >{\raggedright\arraybackslash}p{(\linewidth - 4\tabcolsep) * \real{0.3333}}
  >{\raggedright\arraybackslash}p{(\linewidth - 4\tabcolsep) * \real{0.3333}}
  >{\raggedright\arraybackslash}p{(\linewidth - 4\tabcolsep) * \real{0.3333}}@{}}
\toprule\noalign{}
\begin{minipage}[b]{\linewidth}\raggedright
\end{minipage} & \begin{minipage}[b]{\linewidth}\raggedright
Exponential
\end{minipage} & \begin{minipage}[b]{\linewidth}\raggedright
Power
\end{minipage} \\
\midrule\noalign{}
\endhead
\bottomrule\noalign{}
\endlastfoot
Equation & \(f(t) = A\,e^{-\lambda t}\) & \(f(t) = A\,(1 + t)^{-d}\) \\
Parameters & \(A\): initial level; \(\lambda\): decay rate & \(A\):
initial level; \(d\): decay exponent \\
Linearizing plot & \(\ln f\) vs.~\(t\) is a straight line with slope
\(-\lambda\) & \(\ln f\) vs.~\(\ln(1+t)\) is a straight line with slope
\(-d\) \\
Relative loss rate \(-f'(t)/f(t)\) & \(\lambda\) (constant) &
\(d/(1+t)\) (decreasing with age) \\
Psychological reading & A memory loses the same \emph{proportion} of its
strength per unit time, whatever its age & Older memories lose strength
more slowly; they become more resistant \\
\end{longtable}
}

The \textbf{relative loss rate} (in survival analysis, the
\emph{hazard}) is the key interpretive difference. It is why the choice
between the two functions is a theoretical question, not a curve-fitting
detail.

\textbf{Noise model.} Recall counts are binomial:

\[
y_i \sim \mathrm{Binomial}\big(n, f(t_i; \boldsymbol\theta)\big).
\]

With \(n = 50\) items and \(f = 0.5\), the SD of the observed proportion
is \(\sqrt{0.5 \cdot 0.5 / 50} \approx 0.07\). That is large compared
with the differences between models over short intervals, which is a
first hint that \textbf{design} (which intervals, how many items)
determines whether data can discriminate between models.

\begin{quote}
\textbf{Advanced: the ``+1'' in the power model is not innocent.}
\(A\,t^{-d}\) is infinite at \(t = 0\), so a shift is added. But
\((1 + t)\) depends on the \textbf{unit} of \(t\): measuring in hours
instead of days changes the shape of the curve, not just the scale of
the axis. Some authors therefore write \(A\,(1 + bt)^{-d}\) with a free
scale parameter \(b\), which adds a third parameter and raises the model
comparison questions of Weeks 6--7. Keep the unit fixed (days)
throughout this course.
\end{quote}

\paragraph{Reading the math}

\begin{itemize}
\tightlist
\item
  \(f(t) = A e^{-\lambda t}\): \(A\) is the recall level right after
  study (\(t = 0\)), and \(\lambda\) says how fast it falls. With
  \(\lambda = 0.1\) per day, recall is multiplied by
  \(e^{-0.1} \approx 0.905\) every day: about 9.5\% of what is left is
  lost each day, on day 1 and on day 100 alike.
\item
  \(f(t) = A(1+t)^{-d}\): going from day 0 to day 1 multiplies recall by
  \(2^{-d}\); going from day 9 to day 10 multiplies it only by
  \((10/11)^{d}\), close to 1. Old memories hardly change any more.
\item
  \textbf{Relative loss rate} \(-f'(t)/f(t)\): ``what fraction of the
  remaining memory is lost per day \emph{right now}''. Constant
  (\(\lambda\)) for the exponential, shrinking (\(d/(1+t)\)) for the
  power function. This single number is the psychological difference
  between the models.
\item
  \(y_i \sim \mathrm{Binomial}(n, f(t_i))\): at delay \(t_i\) the
  participant recalls \(y_i\) of \(n\) items, and each item is recalled
  with probability \(f(t_i)\). The observed proportion \(y_i/n\)
  scatters around \(f(t_i)\) with SD \(\sqrt{f(1-f)/n}\), about 0.07 for
  \(n = 50\).
\end{itemize}

\paragraph{For an engineer}

\begin{itemize}
\tightlist
\item
  The \textbf{exponential} curve is the step response of a first-order
  system (an RC circuit discharging): one time constant
  \(\tau = 1/\lambda\). A constant relative loss rate is the
  ``memoryless'' property you know from exponential lifetimes in
  reliability theory.
\item
  The \textbf{power} curve has no single time constant. A useful
  picture: a sum of many exponentials with \emph{different} time
  constants (fast and slow components) looks like a power law. Week 5
  shows this exactly, and it is one candidate explanation of power-law
  forgetting.
\item
  On a \textbf{log--log} plot a power law is a straight line; on a
  \textbf{semi-log} plot an exponential is a straight line. Engineers
  use the same trick to identify system orders and \(1/f\) noise.
\item
  The binomial noise model is the counting statistics of a detector:
  like counting bit errors out of \(n\) bits, the uncertainty of a
  proportion shrinks with \(1/\sqrt{n}\).
\end{itemize}

\subsubsection{4. Simulation as a thinking tool}

Before fitting anything, simulate. Simulation answers questions such as:

\begin{itemize}
\tightlist
\item
  \textbf{What can the model produce, and what can it \emph{not}
  produce?} Equal-variance SDT with \(d' > 0\) can never produce
  \(H < F\), and it can never produce a zROC slope different from 1.
  Both exponential and power curves are always decreasing, so they
  cannot capture \emph{reminiscence} (recall improving over time). The
  set of things a model \textbf{cannot} do is what makes it testable.
\item
  \textbf{Which parameters matter where?} \(A\) moves the whole
  forgetting curve; \(\lambda\) or \(d\) controls the shape.
\item
  \textbf{Can the design discriminate between models?} If two models
  make nearly identical predictions at the intervals we plan to test, no
  fitting method can tell them apart (lab Exercise 4; Assignment A1).
\end{itemize}

\subsubsection{5. Preview: estimation by linearization, and why it is
not enough}

A traditional way to fit forgetting curves is to log-transform the
observed proportions and fit a straight line by ordinary regression. It
is fast and intuitive, and we try it in the lab. It has three problems
that motivate Weeks 3--4:

\begin{enumerate}
\def\labelenumi{\arabic{enumi}.}
\tightlist
\item
  \(\ln(0)\) is undefined, so an interval with zero recalls breaks the
  method.
\item
  The log transform distorts the noise: the scatter of \(\ln \hat p\) is
  much larger for small proportions, but regression treats all points as
  equally reliable.
\item
  The line is fit on the log scale, but we care about errors on the
  original (proportion) scale.
\end{enumerate}

The principled alternative is to fit the model directly to the
untransformed data, by minimizing squared error (Week 3) or maximizing
the binomial likelihood (Week 4).

\paragraph{For the instructor}

\begin{itemize}
\tightlist
\item
  \textbf{Open with the radiologist puzzle} (Motivation) and let
  students vote on who is better before showing any math. Most will pick
  by percent correct, then realize it is identical.
\item
  \textbf{Draw, don't derive.} The two overlapping bell curves with a
  movable vertical line explain \(H\), \(F\), \(d'\) and \(c\) better
  than any formula. Move the line left and right on the board and ask
  what happens to \(H\) and \(F\); then pull the curves apart and ask
  again.
\item
  \textbf{Use \(\Phi\) tables or Python live.} Compute R1's \(d'\) with
  \nolinkurl{norm.ppf} in front of the class. Students who are new to
  statistics need to see that \(\Phi^{-1}\) is ``just a lookup''.
\item
  \textbf{Connect to something they know:} a COVID or pregnancy test has
  a hit rate (sensitivity) and a false-alarm rate (1 − specificity).
  Medical ``sensitivity'' is the hit rate, \emph{not} \(d'\): point out
  this naming clash.
\item
  \textbf{Typical questions.}

  \begin{itemize}
  \tightlist
  \item
    \emph{``Why assume normal distributions?''} It is a convenient and
    often adequate assumption; the zROC test (slope 1) checks it, and
    the unequal-variance model relaxes it.
  \item
    \emph{``Is a liberal criterion bad?''} No: the best criterion
    depends on costs and base rates.
  \item
    \emph{``Exponential or power, which is right?''} For many datasets
    the power function fits better, but averaging over people can create
    a power shape (Week 5). Keep the question open.
  \end{itemize}
\end{itemize}

\begin{center}\rule{0.5\linewidth}{0.5pt}\end{center}

\subsection{Worked example}

\subsubsection{Part 1: Two radiologists (and a third)}

Each radiologist reads 100 tumor scans and 100 healthy scans.

{\def\LTcaptype{none} % do not increment counter
\begin{landscape}\small
\begin{longtable}[]{@{}
  >{\raggedright\arraybackslash}p{(\linewidth - 14\tabcolsep) * \real{0.1250}}
  >{\raggedright\arraybackslash}p{(\linewidth - 14\tabcolsep) * \real{0.1250}}
  >{\raggedright\arraybackslash}p{(\linewidth - 14\tabcolsep) * \real{0.1250}}
  >{\raggedright\arraybackslash}p{(\linewidth - 14\tabcolsep) * \real{0.1250}}
  >{\raggedright\arraybackslash}p{(\linewidth - 14\tabcolsep) * \real{0.1250}}
  >{\raggedright\arraybackslash}p{(\linewidth - 14\tabcolsep) * \real{0.1250}}
  >{\raggedright\arraybackslash}p{(\linewidth - 14\tabcolsep) * \real{0.1250}}
  >{\raggedright\arraybackslash}p{(\linewidth - 14\tabcolsep) * \real{0.1250}}@{}}
\toprule\noalign{}
\begin{minipage}[b]{\linewidth}\raggedright
Observer
\end{minipage} & \begin{minipage}[b]{\linewidth}\raggedright
\(H\)
\end{minipage} & \begin{minipage}[b]{\linewidth}\raggedright
\(F\)
\end{minipage} & \begin{minipage}[b]{\linewidth}\raggedright
\(\Phi^{-1}(H)\)
\end{minipage} & \begin{minipage}[b]{\linewidth}\raggedright
\(\Phi^{-1}(F)\)
\end{minipage} & \begin{minipage}[b]{\linewidth}\raggedright
\(d'\)
\end{minipage} & \begin{minipage}[b]{\linewidth}\raggedright
\(c\)
\end{minipage} & \begin{minipage}[b]{\linewidth}\raggedright
Percent correct
\end{minipage} \\
\midrule\noalign{}
\endhead
\bottomrule\noalign{}
\endlastfoot
R0 (neutral) & 0.84 & 0.16 & 0.994 & −0.994 & 1.99 & 0.00 & 84\% \\
R1 (liberal) & 0.93 & 0.31 & 1.476 & −0.496 & 1.97 & −0.49 & 81\% \\
R2 (conservative) & 0.69 & 0.07 & 0.496 & −1.476 & 1.97 & +0.49 &
81\% \\
\end{longtable}
\end{landscape}
}

\textbf{Step by step for R1:} 1. \(\Phi^{-1}(0.93) = 1.476\) and
\(\Phi^{-1}(0.31) = -0.496\) (\nolinkurl{scipy.stats.norm.ppf}). 2.
\(d' = 1.476 - (-0.496) = 1.972\). 3.
\(c = -\tfrac{1}{2}(1.476 + (-0.496)) = -0.490\), which is negative, so
R1 is liberal. 4. Check by plugging back in:
\(H = \Phi(1.972/2 + 0.490) = \Phi(1.476) = 0.93\). \ensuremath{\checkmark}

\textbf{Interpretation.} All three radiologists have essentially the
same sensitivity (\(d' \approx 2\)). R1 and R2 have the \emph{same}
percent correct but \emph{opposite} biases. Which bias is better is not
a question of perception but of costs: missing a tumor is usually far
worse than a false alarm, so a liberal criterion may be the
\emph{desirable} one. SDT lets us ask ``can they see it?'' and ``how
willing are they to say so?'' separately.

\textbf{The extreme-rate fix.} A participant scores 50/50 hits and 3/50
false alarms. Then \(H = 1\) and \(d' = \infty\). With the log-linear
correction: \(H = 50.5/51 = 0.990\), \(F = 3.5/51 = 0.069\), so
\(d' = 2.33 - (-1.49) = 3.82\).

\subsubsection{Part 2: Exponential vs.~power forgetting}

Take \(A = 0.95\) for both models, \(\lambda = 0.25\)/day for the
exponential and \(d = 0.6\) for the power model.

{\def\LTcaptype{none} % do not increment counter
\begin{landscape}\small
\begin{longtable}[]{@{}
  >{\raggedright\arraybackslash}p{(\linewidth - 12\tabcolsep) * \real{0.1429}}
  >{\raggedright\arraybackslash}p{(\linewidth - 12\tabcolsep) * \real{0.1429}}
  >{\raggedright\arraybackslash}p{(\linewidth - 12\tabcolsep) * \real{0.1429}}
  >{\raggedright\arraybackslash}p{(\linewidth - 12\tabcolsep) * \real{0.1429}}
  >{\raggedright\arraybackslash}p{(\linewidth - 12\tabcolsep) * \real{0.1429}}
  >{\raggedright\arraybackslash}p{(\linewidth - 12\tabcolsep) * \real{0.1429}}
  >{\raggedright\arraybackslash}p{(\linewidth - 12\tabcolsep) * \real{0.1429}}@{}}
\toprule\noalign{}
\begin{minipage}[b]{\linewidth}\raggedright
\(t\) (days)
\end{minipage} & \begin{minipage}[b]{\linewidth}\raggedright
0
\end{minipage} & \begin{minipage}[b]{\linewidth}\raggedright
1
\end{minipage} & \begin{minipage}[b]{\linewidth}\raggedright
2
\end{minipage} & \begin{minipage}[b]{\linewidth}\raggedright
7
\end{minipage} & \begin{minipage}[b]{\linewidth}\raggedright
14
\end{minipage} & \begin{minipage}[b]{\linewidth}\raggedright
30
\end{minipage} \\
\midrule\noalign{}
\endhead
\bottomrule\noalign{}
\endlastfoot
Exponential \(0.95\,e^{-0.25t}\) & 0.950 & 0.740 & 0.576 & 0.165 & 0.029
& 0.001 \\
Power \(0.95\,(1+t)^{-0.6}\) & 0.950 & 0.627 & 0.491 & 0.273 & 0.187 &
0.121 \\
\end{longtable}
\end{landscape}
}

\textbf{Half-life (time until recall halves from its initial level):} -
Exponential:
\(e^{-\lambda t} = 1/2 \Rightarrow t_{1/2} = \ln 2 / \lambda = 0.693/0.25 = 2.77\)
days. - Power:
\((1+t)^{-d} = 1/2 \Rightarrow t_{1/2} = 2^{1/d} - 1 = 2^{1.667} - 1 = 2.17\)
days.

\textbf{Interpretation.} Over the first two days the curves are fairly
similar: they differ by about 0.1, which is only about 1.5 binomial SDs
with \(n = 50\) items. After a month they disagree completely: the
exponential model says practically everything is forgotten (0.1\%),
while the power model says 12\% is still retained. \textbf{Long
retention intervals are where these models can be told apart}, a design
lesson you will use in Assignment A1.

\begin{center}\rule{0.5\linewidth}{0.5pt}\end{center}

\subsection{Lab (2 hours)}

\subsubsection{Goals}

\begin{enumerate}
\def\labelenumi{\arabic{enumi}.}
\tightlist
\item
  Implement SDT as a generative model, simulate yes/no experiments, and
  estimate \(d'\) and \(c\).
\item
  Use ROC and zROC plots to \textbf{evaluate} the equal-variance
  assumption.
\item
  Implement exponential and power forgetting with binomial noise, and
  explore their parameters.
\item
  \textbf{Fit} both forgetting models by linearization,
  \textbf{evaluate} them with residuals, and \textbf{interpret} the
  result.
\end{enumerate}

\textbf{Timing:} Part A 30 min · Part B 25 min · Part C 30 min · Part D
25 min · remaining time: exercises.

\subsubsection{Setup}

\begin{quote}
\textbf{Code · \RL{\CMfa کد:}} \nolinkurl{code/week02.py}, cell 1 · \RL{\CMfa سلول}
1
\end{quote}

\subsubsection{Part A: SDT as a generative model (30 min)}

\begin{quote}
\textbf{Code · \RL{\CMfa کد:}} \nolinkurl{code/week02.py}, cell 2 · \RL{\CMfa سلول}
2
\end{quote}

Simulate one participant and recover the parameters:

\begin{quote}
\textbf{Code · \RL{\CMfa کد:}} \nolinkurl{code/week02.py}, cell 3 · \RL{\CMfa سلول}
3
\end{quote}

\textbf{Expected output:}

\Needspace{6\baselineskip}
\begin{Verbatim}
Counts: {'hits': 68, 'misses': 32, 'fas': 11, 'crs': 89}
Predicted H = 0.674, F = 0.147
Estimated d' = 1.67 (true 1.5),  c = 0.37 (true 0.3)
\end{Verbatim}

One simulated participant tells us little about how good the
\emph{estimator} is. Let's repeat the experiment many times: this is a
preview of the \textbf{parameter recovery} logic of Week 7.

\begin{quote}
\textbf{Code · \RL{\CMfa کد:}} \nolinkurl{code/week02.py}, cell 4 · \RL{\CMfa سلول}
4
\end{quote}

\textbf{Expected output:}

\Needspace{9\baselineskip}
\begin{Verbatim}
            d_hat         c_hat
             mean    std   mean    std
n_per_type
20          1.500  0.449  0.296  0.224
50          1.489  0.290  0.295  0.143
200         1.502  0.142  0.299  0.071
\end{Verbatim}

\textbf{What to notice:} the estimates scatter around the true values,
and the scatter shrinks as the number of trials grows. With 20 trials
per type, the SD of \(\hat d'\) is about 0.45, so individual estimates
are often off by \ensuremath{\pm}0.5 or more, and roughly 4× more trials are needed to
halve the scatter. On average the estimates are on target here. With
higher \(d'\) and few trials, however, extreme rates become common, and
the correction then pulls estimates toward 0 (try \nolinkurl{true_d = 3}
with 20 trials). Keep this in mind whenever you compare \(d'\) across
groups tested with different numbers of trials.

\subsubsection{Part B: ROC curves and the equal-variance assumption (25
min)}

First, the model's predicted ROCs:

\begin{quote}
\textbf{Code · \RL{\CMfa کد:}} \nolinkurl{code/week02.py}, cell 5 · \RL{\CMfa سلول}
5
\end{quote}

Now simulate a \textbf{confidence-rating} experiment. Responses are
given on a 6-point scale (1 = ``sure noise'' \ldots{} 6 = ``sure
signal''), which amounts to five criteria at once. We generate data from
an equal-variance observer and from an \textbf{unequal-variance}
observer whose signal distribution has SD 1.5.

\begin{quote}
\textbf{Code · \RL{\CMfa کد:}} \nolinkurl{code/week02.py}, cell 6 · \RL{\CMfa سلول}
6
\end{quote}

\textbf{Expected output:}

\Needspace{5\baselineskip}
\begin{Verbatim}
sigma_signal = 1.0: zROC slope = 0.96 (predicted 1.00); d' at each criterion = [1.52 1.48 1.48 1.43 1.43]
sigma_signal = 1.5: zROC slope = 0.67 (predicted 0.67); d' at each criterion = [1.46 1.34 1.16 0.96 0.84]
\end{Verbatim}

\textbf{Evaluate and interpret.} For the equal-variance observer, the
zROC slope is close to 1 and \(d'\) is about the same at every
criterion: the model's assumption holds. For the unequal-variance
observer, the slope is close to \(1/1.5 = 0.67\) and \(d'\)
\textbf{changes systematically with the criterion}. If we had collected
only yes/no data from such an observer, a group difference in ``\(d'\)''
could be produced purely by a difference in bias. Checking a model's
signature prediction (here: zROC slope = 1) is a basic form of
\textbf{model evaluation}.

\subsubsection{Part C: Forgetting curves with binomial noise (30 min)}

\begin{quote}
\textbf{Code · \RL{\CMfa کد:}} \nolinkurl{code/week02.py}, cell 7 · \RL{\CMfa سلول}
7
\end{quote}

Explore the parameters: how do \(A\) and the decay parameter change each
curve?

\begin{quote}
\textbf{Code · \RL{\CMfa کد:}} \nolinkurl{code/week02.py}, cell 8 · \RL{\CMfa سلول}
8
\end{quote}

Now generate a \textbf{simulated experiment}. The data come from the
\emph{power} model, with 50 items tested at each of 8 retention
intervals:

\begin{quote}
\textbf{Code · \RL{\CMfa کد:}} \nolinkurl{code/week02.py}, cell 9 · \RL{\CMfa سلول}
9
\end{quote}

\textbf{Expected output:}

\Needspace{12\baselineskip}
\begin{Verbatim}
       t  recalled   n  p_obs
0   0.25        42  50   0.84
1   1.00        30  50   0.60
2   2.00        26  50   0.52
3   4.00        25  50   0.50
4   7.00        13  50   0.26
5  14.00        13  50   0.26
6  21.00        13  50   0.26
7  30.00        14  50   0.28
\end{Verbatim}

\subsubsection{Part D: Fit by linearization, evaluate, interpret (25
min)}

\textbf{Fit.} Fit each model by a straight-line regression on the log
scale (Lecture §5).

\begin{quote}
\textbf{Code · \RL{\CMfa کد:}} \nolinkurl{code/week02.py}, cell 10 · \RL{\CMfa سلول}
10
\end{quote}

\textbf{Evaluate.} Look at the fitted curves \emph{and} the residuals. A
good model leaves residuals that scatter randomly around zero; a
systematic pattern means the model's \emph{shape} is wrong.

\begin{quote}
\textbf{Code · \RL{\CMfa کد:}} \nolinkurl{code/week02.py}, cell 11 · \RL{\CMfa سلول}
11
\end{quote}

\textbf{Expected output:}

\Needspace{6\baselineskip}
\begin{Verbatim}
exp: A = 0.550, rate = 0.032, RMSD = 0.1308
pow: A = 0.797, rate = 0.370, RMSD = 0.0628
true (power): A = 0.9, d = 0.5
\end{Verbatim}

\textbf{Interpret (write 2--3 sentences for each in your notebook):}

\begin{enumerate}
\def\labelenumi{\arabic{enumi}.}
\tightlist
\item
  Which model describes the data better, and what does the residual
  pattern of the \emph{other} model look like? (Look for an inverted-U
  or U shape: the exponential overshoots at intermediate delays and
  undershoots at the long ones, or the reverse.)
\item
  Using the fitted parameters, what is each model's relative loss rate
  after 1 day and after 20 days? What does this mean for memory?
\item
  The power model is the true model here, yet its estimates
  (\(\hat A \approx 0.80\), \(\hat d \approx 0.37\)) are noticeably off
  the true values (\(A^\ast = 0.9\), \(d^\ast = 0.5\)). Name two
  reasons. (Hints: look at the last three data points; think about what
  the log transform does to the weights of small proportions.) In Week 3
  we refit these same data on the original scale and compare.
\end{enumerate}

\subsubsection{Exercises (increasing difficulty)}

\textbf{Exercise 1 \ensuremath{\star} (extreme rates).} A participant has 40/40 hits and
0/40 false alarms. Compute \(d'\) without and with the log-linear
correction. What does the corrected \(d'\) represent, and why should you
report the correction you used?

\textbf{Exercise 2 \ensuremath{\star} (half-lives).} Using your fitted parameters from
Part D, compute each model's half-life. Then compute how long it takes
recall to drop to 10\% of its initial level. Where do the models
disagree most?

\textbf{Exercise 3 \ensuremath{\star}\ensuremath{\star} (Ebbinghaus's data).} The classic savings data
from Ebbinghaus's 1885 self-experiments are often reported as follows
{[}verify against the original source before using them in class{]}:

{\def\LTcaptype{none} % do not increment counter
\begin{landscape}\small
\begin{longtable}[]{@{}llllllll@{}}
\toprule\noalign{}
Interval & 20 min & 1 h & 8.8 h & 1 day & 2 days & 6 days & 31 days \\
\midrule\noalign{}
\endhead
\bottomrule\noalign{}
\endlastfoot
Savings (\%) & 58.2 & 44.2 & 35.8 & 33.7 & 27.8 & 25.4 & 21.1 \\
\end{longtable}
\end{landscape}
}

Convert the intervals to days, fit both models by linearization, and
compare their residuals. Which shape fits better? Then redo the analysis
with time measured in \textbf{hours}. Does the power fit change?
Explain, using the Advanced box in Lecture §3.

\textbf{Exercise 4 \ensuremath{\star}\ensuremath{\star} (which design can tell the models apart?).} Choose
\((A, \lambda)\) and \((A, d)\) so that the two models predict the
\emph{same} recall at \(t = 1\) and \(t = 7\) days. (Hint: two equations
in two unknowns per model; take logs.) Then plot both curves from
\(t = 0\) to \(60\) days and mark where their difference exceeds 2
binomial SDs with \(n = 50\). Which retention intervals would you test?

\textbf{Exercise 5 \ensuremath{\star}\ensuremath{\star}\ensuremath{\star} (Advanced: zero counts).} Generate data from the
exponential model with \(A = 0.9\), \(\lambda = 0.3\) at the intervals
in Part C, with \(n = 20\) items. How often (out of 1000 simulated
datasets) does at least one interval produce zero recalls, which breaks
\nolinkurl{fit_linearized}? Propose and implement a quick fix (e.g.,
adding 0.5 to counts) and check whether it biases the estimate of
\(\lambda\). Explain why this is a sign that we need a method that works
on the original scale (Weeks 3--4).

\subsubsection{Lab check-off}

Show the instructor or TA: (1) your SDT recovery table from Part A, (2)
the two empirical zROCs with their slopes from Part B, and (3) the
residual plot from Part D with a one-sentence interpretation.

\begin{center}\rule{0.5\linewidth}{0.5pt}\end{center}

\subsection{Connection to AI \& LLMs}

\textbf{AI-text detectors are observers in an SDT task.} Universities
increasingly use tools that score a text for ``how likely it was written
by AI''. Each text is either AI-written (``signal'') or human-written
(``noise''), and the tool says ``AI'' when its score exceeds a
threshold. This is exactly the SDT setting of this week: the tool has a
\textbf{sensitivity} \(d'\) (how well the score separates the two kinds
of texts) and a \textbf{criterion} \(c\) (the threshold the user
chooses). A false alarm means a student is wrongly accused, so the
choice of criterion is an ethical decision, not only a technical one.

\begin{quote}
\textbf{Code · \RL{\CMfa کد:}} \nolinkurl{code/week02.py}, cell 12 · \RL{\CMfa سلول}
12
\end{quote}

\textbf{Expected output:}

\Needspace{6\baselineskip}
\begin{Verbatim}
threshold 0.0: H = 0.84, F = 0.49, d' = 1.03, c = -0.50
threshold 0.6: H = 0.70, F = 0.28, d' = 1.11, c = 0.03
threshold 1.2: H = 0.51, F = 0.09, d' = 1.35, c = 0.66
\end{Verbatim}

Moving the threshold changes \(c\) a lot (from −0.50 to 0.66), while
\(d'\) stays near its true value of 1.2. The variation in \(d'\) is
sampling noise from 200 texts per group.

\textbf{Try it (no API needed).} - \ensuremath{\star} Suppose 10\% of submitted essays
are AI-written. At threshold 1.2, what is
\(P(\text{AI-written} \mid \text{flagged})\)? (Use Bayes' rule with
\(H\) and \(F\).) Would you act on a flag? - \ensuremath{\star}\ensuremath{\star} Plot the ROC of the
detector (all thresholds from −2 to 3) and its zROC. Is the slope close
to 1? Repeat with human-text scores of SD 1.5 (an unequal-variance
detector). - \ensuremath{\star}\ensuremath{\star} Discussion: detectors have been reported to flag texts
by non-native English writers more often {[}verify: Liang et al., 2023,
\emph{Patterns}{]}. In SDT terms, what would that mean, and what should
a university do?

\begin{center}\rule{0.5\linewidth}{0.5pt}\end{center}

\subsection{Common mistakes \& misconceptions}

\begin{enumerate}
\def\labelenumi{\arabic{enumi}.}
\tightlist
\item
  \textbf{``Percent correct is good enough.''} It confounds sensitivity
  and bias (see the radiologists). Use \(d'\) and \(c\) whenever
  responses can be biased.
\item
  \textbf{Computing \(d'\) as \(\Phi^{-1}(H - F)\) or from percentages.}
  The formula is \(\Phi^{-1}(H) - \Phi^{-1}(F)\), with rates as
  proportions between 0 and 1.
\item
  \textbf{Forgetting the extreme-rate problem.} \(H = 1\) or \(F = 0\)
  gives an infinite \(d'\); Python returns \nolinkurl{inf} without an
  error. Always use (and report) a correction.
\item
  \textbf{Assuming \(d'\) is bias-free no matter what.} It is
  criterion-independent only under the equal-variance model. If the zROC
  slope differs from 1, \(d'\) from one \((H, F)\) pair depends on bias.
\item
  \textbf{``Exponential and power are just two curve shapes.''} They
  make different claims about how forgetting works (constant
  vs.~decreasing relative loss rate).
\item
  \textbf{Judging a model only by its fit statistic.} A small RMSD with
  a systematic residual pattern means the model's shape is wrong. Always
  plot the residuals.
\item
  \textbf{Ignoring the noise model.} A model is not just
  \(f(x;\boldsymbol\theta)\). The binomial noise of recall proportions
  is part of the model, and it tells you which differences are
  meaningful.
\item
  \textbf{Changing the units of \(t\) in the shifted power model} and
  expecting the same shape.
\end{enumerate}

\begin{center}\rule{0.5\linewidth}{0.5pt}\end{center}

\subsection{Check your understanding}

\textbf{1.} Name the five components of a model and give the SDT version
of each. \textbf{Answer} Design (signal vs.~noise trials; yes/no
responses); parameters (free: \(d'\), \(c\); fixed: noise mean 0, both
SDs 1); deterministic core (\(H = \Phi(d'/2 - c)\),
\(F = \Phi(-d'/2 - c)\)); noise model (hit and FA counts are binomial);
linking assumption (respond ``yes'' when the evidence exceeds the
criterion).

\textbf{2.} Why are the noise mean and the SDs in SDT fixed rather than
free? \textbf{Answer} Because they are not identifiable: shifting both
distributions and the criterion by the same amount, or stretching the
whole axis, leaves every prediction unchanged. Fixing them (scaling
constraints) defines the units in which \(d'\) and \(c\) are measured.

\textbf{3.} An observer has \(H = 0.80\) and \(F = 0.20\). Compute
\(d'\) and \(c\). (\(\Phi^{-1}(0.80) = 0.842\)) \textbf{Answer}
\(d' = 0.842 - (-0.842) = 1.68\);
\(c = -\tfrac{1}{2}(0.842 - 0.842) = 0\), so the observer is unbiased.

\textbf{4.} What does equal-variance SDT predict about the zROC, and
what does it mean if the empirical slope is 0.7? \textbf{Answer} A
straight line with slope 1. A slope of 0.7 suggests the signal
distribution is more variable than the noise distribution
(\(\sigma_s \approx 1/0.7 \approx 1.4\)). Then single-point \(d'\)
depends on the criterion, and an unequal-variance model (or \(d_a\))
should be used.

\textbf{5.} What is the relative loss rate of each forgetting model, and
why does it matter psychologically? \textbf{Answer} Exponential:
\(\lambda\), constant. Power: \(d/(1+t)\), decreasing over time. The
exponential says memories lose the same proportion of strength per day
at every age; the power function says older memories become more
resistant to further forgetting.

\textbf{6.} Which plot would you make to check quickly whether
forgetting data look exponential? And for power? \textbf{Answer}
Exponential: \(\ln(\text{recall})\) against \(t\) (semi-log), which
should be a straight line. Power: \(\ln(\text{recall})\) against
\(\ln(1+t)\) (log-log), which should be a straight line.

\textbf{7.} Give two reasons why fitting a straight line to
log-transformed proportions is a poor estimation method. \textbf{Answer}
(i) It breaks when an observed proportion is 0. (ii) It distorts the
noise: small proportions become very variable on the log scale but are
weighted equally, and the errors being minimized are on the wrong (log)
scale. Fitting on the original scale with least squares or maximum
likelihood avoids these problems.

\begin{center}\rule{0.5\linewidth}{0.5pt}\end{center}

\subsection{Assignment}

\subsubsection{A1: From verbal to formal (1.0 point of 20; released Week
2, due at the start of the Week 4 lecture)}

\textbf{Scenario.} Two research groups study forgetting of vocabulary
learned in a foreign-language course. - \textbf{Group C} says:
\emph{``Memories decay at a constant rate.''} - \textbf{Group R} says:
\emph{``Memories become more resistant to forgetting as they age.''}

Both groups claim that the published data (recall drops quickly and then
levels off) support their view.

\textbf{Tasks.} Submit a Jupyter notebook plus a PDF export.

\begin{enumerate}
\def\labelenumi{\arabic{enumi}.}
\tightlist
\item
  \textbf{Formalize (Group C and Group R).} Write each verbal claim as a
  formal model with all five components from Lecture §1: design,
  free/fixed parameters with ranges, deterministic core, binomial noise
  model, and linking assumption. Justify in 2--3 sentences why each
  equation captures its verbal claim (use the relative loss rate).
\item
  \textbf{Implement.} Write Python functions that (a) return predicted
  recall probabilities and (b) simulate binomial recall counts for given
  intervals and number of items \(n\).
\item
  \textbf{Calibrate.} Both groups agree that recall is \textbf{0.60
  after 1 day} and \textbf{0.35 after 7 days}. Derive, by hand and in
  closed form, the parameters of each model that reproduce these two
  values exactly. Show the derivation (take logs) and verify it
  numerically.
\item
  \textbf{Find where they diverge.} With the calibrated parameters, plot
  both predicted curves from 0 to 90 days. For \(n = 40\) items, compute
  at each interval a \emph{discriminability index}
  \(D(t) = \big(f_C(t) - f_R(t)\big)^2 / \mathrm{Var}(t)\), where
  \(\mathrm{Var}(t) = \bar f(1 - \bar f)/n\) and \(\bar f\) is the
  average of the two predictions. Plot \(D(t)\).
\item
  \textbf{Design an experiment.} You can test each participant at only
  \textbf{4} retention intervals, chosen from \{0.5, 1, 2, 4, 7, 14, 30,
  60, 90\} days. Choose the set that maximizes the summed \(D(t)\),
  explain why it makes sense, and simulate one dataset from \emph{each}
  model at your design. Show the two datasets on one plot.
\item
  \textbf{Interpret (max. 300 words).} Which data would favor Group C
  and which Group R? What does your analysis say about why both groups
  can claim support from ``the published data''? Name one assumption of
  your models that could be wrong, and how it would change your
  conclusions.
\end{enumerate}

\textbf{Grading rubric (10 points, scaled to 1.0 of the final grade):}

{\def\LTcaptype{none} % do not increment counter
\begin{longtable}[]{@{}
  >{\raggedright\arraybackslash}p{(\linewidth - 4\tabcolsep) * \real{0.3333}}
  >{\raggedright\arraybackslash}p{(\linewidth - 4\tabcolsep) * \real{0.3333}}
  >{\raggedright\arraybackslash}p{(\linewidth - 4\tabcolsep) * \real{0.3333}}@{}}
\toprule\noalign{}
\begin{minipage}[b]{\linewidth}\raggedright
Criterion
\end{minipage} & \begin{minipage}[b]{\linewidth}\raggedright
Points
\end{minipage} & \begin{minipage}[b]{\linewidth}\raggedright
Full marks require
\end{minipage} \\
\midrule\noalign{}
\endhead
\bottomrule\noalign{}
\endlastfoot
Formalization (Task 1) & 2 & All five components for both models;
correct parameter ranges; the link between equation and verbal claim is
explained via the relative loss rate \\
Code (Task 2) & 2 & Correct, commented, reproducible (fixed seeds);
functions are vectorized or clearly looped \\
Calibration (Task 3) & 2 & Correct closed-form derivation and numerical
check for both models \\
Divergence and design (Tasks 4--5) & 2.5 & Correct \(D(t)\); the chosen
design is justified and actually maximizes the criterion; clear plots \\
Interpretation (Task 6) & 1.5 & Clear, specific, refers to your own
results; names a real limitation \\
\end{longtable}
}

Late submissions lose 10\% per day (max. 3 days).

\begin{center}\rule{0.5\linewidth}{0.5pt}\end{center}

\subsection{Further reading}

\begin{itemize}
\tightlist
\item
  \textbf{Farrell \& Lewandowsky (2018), chapter 2.} Building a model
  from a verbal theory; the components of a model and the role of
  parameters. \emph{(Main reading.)}
\item
  \textbf{Lee \& Wagenmakers (2014), chapter 11.} Signal detection
  theory treated as a (Bayesian) cognitive model; we return to it in
  Weeks 9 and 13. Read the introductory part now, skip the code.
\item
  \textbf{Verguts (2022)} and \textbf{Anderson (2014)}: the chapters on
  model building and parameters {[}verify chapter numbers{]}.
\item
  \emph{Background on forgetting curves:} Lee \& Wagenmakers (2014),
  chapter 10 (memory retention) uses a closely related exponential
  retention model; we will fit it in a Bayesian way in Weeks 9--11.
\end{itemize}

\emph{Preview for Week 3:} we drop linearization and fit models directly
by minimizing squared error, which means learning how optimization
algorithms search a parameter space.
